% !TEX program = pdflatex
% Compile: pdflatex thm44_verification_report.tex (twice for TOC); xelatex/lualatex also work
\documentclass[11pt]{article}

\usepackage[a4paper,margin=1in]{geometry}
\usepackage{amsmath,amssymb,amsthm,mathtools}
\usepackage{microtype}
\usepackage{booktabs}
\usepackage{array}
\usepackage{longtable}
\usepackage{tabularx}
\usepackage{xcolor}
\usepackage{colortbl}
\usepackage{enumitem}
\usepackage[hidelinks,bookmarks=true]{hyperref}
\setlength{\tabcolsep}{4pt}
\renewcommand{\arraystretch}{1.15}

\newtheorem{theorem}{Theorem}

% ------------------------------------------------------------
% Colors and verdict / hypothesis macros
% ------------------------------------------------------------
\definecolor{cConfirm}{HTML}{2E7D32}
\definecolor{cInapp}{HTML}{E65100}
\definecolor{cLift}{HTML}{1565C0}
\definecolor{cCounter}{HTML}{C62828}
\definecolor{cPassBg}{HTML}{E8F5E9}
\definecolor{cFailBg}{HTML}{FFCDD2}
\definecolor{cNABg}{HTML}{ECEFF1}
\definecolor{cSharpBg}{HTML}{FFF8E1}

\newcommand{\vCONF}{\textcolor{cConfirm}{\textsc{Confirms}}}
\newcommand{\vINAP}{\textcolor{cInapp}{\textsc{Inappl.}}}
\newcommand{\vLIFT}{\textcolor{cLift}{\textsc{Not-Liftable}}}
\newcommand{\vCNTR}{\textcolor{cCounter}{\textsc{Counterex.}}}

\newcommand{\yes}{\textcolor{cConfirm}{$\checkmark$}}
\newcommand{\no}{\textcolor{cCounter}{$\boldsymbol{\times}$}}
\newcommand{\rOK}{\cellcolor{cPassBg}\yes}
\newcommand{\rFAIL}{\cellcolor{cFailBg}\no}
\newcommand{\rNA}{\cellcolor{cNABg}\,--\,}
\newcommand{\enot}[2]{\ensuremath{#1\!\times\!10^{#2}}}

\newcolumntype{C}[1]{>{\centering\arraybackslash}p{#1}}
\newcolumntype{L}[1]{>{\raggedright\arraybackslash}p{#1}}

\title{\bfseries Verification of the Two--Anchor Bell--Flag Ladder Theorem\\
  on the Explicit Codes of \texttt{arXiv:2504.20847}}
\author{Numerical report (43 codes, 5 rungs, tolerance $10^{-8}$)\\
  generated by \texttt{verify\_thm44\_paper.py}}
\date{}

\begin{document}
\maketitle

\begin{abstract}
\noindent
We verify Theorem~4.4 of \texttt{dyadic\_lifts.tex} (the two-anchor Bell-flag
ladder theorem) on every explicit non-additive $K=2$ code whose
computational-basis $|0_L\rangle$ is given in \texttt{arXiv:2504.20847}
(Zhang--Wu--Huang--Zeng~2025). Each code is recursively dyadic Bell-lifted to
rung 5; at every rung all Pauli errors of weight $\le 2$ are enumerated and the
Knill--Laflamme distance-3 condition is checked with numerical tolerance
$10^{-8}$. The headline result is

\smallskip
\centerline{\textbf{30 \vCONF{} \;\; 0 \vCNTR{} \;\; 9 \vINAP{} \;\;
4 \vLIFT{}}}\smallskip

\noindent
Every code for which the theorem makes a claim (30 cases) preserves distance 3
at all 5 rungs at machine precision; no counterexample is found across the
remaining 13 codes either, only failures or scopes consistent with the
theorem's stated assumptions.
\end{abstract}

\tableofcontents

\bigskip
\hrule

% ============================================================
\section{Executive summary}
\label{sec:summary}
% ============================================================

\paragraph{Theorem under test.}
The two-anchor Bell-flag ladder theorem (\texttt{dyadic\_lifts.tex},
Theorem~4.4) asserts: if a seed SS-LP code satisfies five hypotheses
(see~\S\ref{tab:supp:hyps}), every dyadic Bell lift of it preserves distance
$\ge 3$. The construction adds 2 qubits per lift and doubles the modulus,
yielding the ladder $(n_r, m_r) = (n_0+2r,\, 2^r m_0)$.

\paragraph{Universe of codes.}
We test all 43 codes from \texttt{arXiv:2504.20847} whose
computational-basis logical-$0$ codeword is given explicitly in the paper. The
codes range over $n\in\{6,7,8\}$, $m\in\{2,4,5,6,\ldots,42\}$, transversal
groups $\mathrm{BD}_{2m}$, $2T$, $2O$, $2I$, $C_{10}$.

\paragraph{Verification pipeline.}
\texttt{paper\_seeds\_2504\_20847.py} encodes each code as a sparse JSON
dictionary; \texttt{verify\_thm44\_paper.py} performs the recursive Bell lift
\eqref{eq:supp:lift} and the weight-$\le2$ Pauli enumeration. Numerical settings
are \texttt{tol}=$10^{-8}$, FP64, \texttt{max\_rungs}=5.

\paragraph{Bottom line.}
\begin{itemize}[leftmargin=2em]
\item \textbf{30 \vCONF.} All five hypotheses hold \emph{and} every rung
preserves distance 3. Worst KL violation across all 30 codes $\times$ 5 rungs
$=$ \enot{8.0}{-14} (the single outlier
\texttt{bd18\_type\_II}, due to numerically resolved phases); the other 28
codes are all at machine precision $\le$ \enot{6.3}{-16}.
\item \textbf{0 \vCNTR.} \emph{No code is found for which all hypotheses hold
but some rung breaks distance 3.}
\item \textbf{9 \vINAP.} At least one hypothesis fails. 8 of these
also have the conclusion fail at some rung (consistent with the necessity
side of the theorem); 1 is the sharpness exhibit \texttt{2i\_perm\_inv}
(\S\ref{ssec:supp:sharpness}) where the conclusion happens to still hold.
\item \textbf{4 \vLIFT.} Codes that are not SS-LP at all; the dyadic Bell lift
is undefined. Only rung-0 distance-3 is recorded.
\end{itemize}

% ============================================================
%% ---- data section taken from the corrected si.tex ----
\section{Theorem under test, hypotheses, and verdict taxonomy}
\label{sec:supp:taxonomy}
%%====================================================================

\paragraph{Theorem 4.4 (paraphrased; see main text \S 3).}
Let $m_0\ge 2$ and $a^{(0)}=(a_1,\ldots,a_{n_0})\in\mathbb{Z}_{m_0}^{n_0}$.
Let $|0^{(0)}\rangle$ and $|1^{(0)}\rangle$ be normalized codewords supported
on the subset-sum sectors $S_0(a^{(0)};m_0)$ and $S_{m_0-1}(a^{(0)};m_0)$
respectively, and assume
\textbf{(H1)} $\sum_j a_j = 2m_0-1$;
\textbf{(H2)} $0\le a_j\le m_0-2$;
\textbf{(H3)} $0^{n_0}$ is the only support string of $|0^{(0)}\rangle$ with
integer subset sum $0$;
\textbf{(H4)} $1^{n_0}$ is the only support string of $|1^{(0)}\rangle$ with
integer subset sum $2m_0-1$;
\textbf{(H5)} both $|i^{(0)}\rangle$ and the first lifted rung
$|i^{(1)}\rangle$ satisfy the Knill--Laflamme equations for all Pauli errors
of weight $\le 2$.
Then every rung obtained by iterating the dyadic Bell lift
\begin{equation}
\label{eq:supp:lift}
|k_L^{\uparrow}\rangle = \!\!\!\sum_{s\in S_k(a;m)} c_s^{(k)}
|s\rangle\,|\Phi_{h(s)\oplus k}\rangle, \qquad m\mapsto 2m,\ a\mapsto (a,m,m),
\end{equation}
satisfies the weight-$\le 2$ Knill--Laflamme equations; in particular each
rung has distance $\ge 3$.

\paragraph{Hypothesis labels used in this supplement.}
Table~\ref{tab:supp:hyps} records the short labels used throughout. Tests are
performed on the seed only (rung $0$); the rung-$1$ KL pass is part of (H5).

\begin{table}[h]
\centering\small
\caption{Hypothesis labels and tests.}\label{tab:supp:hyps}
\begin{tabular}{@{}clL{7.4cm}@{}}
\toprule
Label & Test & Meaning \\
\midrule
\textbf{strong} & $\sum_j a_j = 2m_0 - 1$        & Strong total-sum identity (H1).\\
\textbf{bnd}    & $\max_j a_j \le m_0 - 2$       & Bounded-entries condition (H2).\\
\textbf{(i)}    & $A_a(s) = 0 \Rightarrow s = 0^{n_0}$ &
                  Unique bottom anchor in $\operatorname{supp}|0^{(0)}\rangle$ (H3).\\
\textbf{(ii)}   & $A_a(s) = 2m_0{-}1 \Rightarrow s = 1^{n_0}$ &
                  Unique top anchor in $\operatorname{supp}|1^{(0)}\rangle$ (H4).\\
\textbf{(iii)}  & rungs $0$ and $1$ pass weight-$\le 2$ KL &
                  Seed and first-lift distance-3 (H5).\\
\bottomrule
\end{tabular}
\end{table}

\paragraph{Verdict taxonomy.}
Table~\ref{tab:supp:verdicts} records the four verdicts we use. Note that
\vINAP{} only states that the theorem makes no claim; we additionally record
whether the conclusion (all rungs $d\ge 3$) happens to hold empirically.

\begin{table}[h]
\centering\small
\caption{Verdict categories used in this supplement.}\label{tab:supp:verdicts}
\begin{tabular}{@{}lL{12.2cm}@{}}
\toprule
Verdict & Meaning \\
\midrule
\vCONF{}  & All five hypotheses hold \emph{and} every rung $r\le 5$ preserves $d\ge 3$. Theorem~4.4 holds on this code.\\
\addlinespace
\vCNTR{}  & All five hypotheses hold \emph{but} some rung $r\le 5$ breaks $d\ge 3$. Would falsify Theorem~4.4.\\
\addlinespace
\vINAP{}  & At least one hypothesis fails $\Rightarrow$ Theorem~4.4 makes no claim. We still record whether the conclusion happens to hold.\\
\addlinespace
\vLIFT{}  & Code is not SS-LP (no diagonal-gate structure) $\Rightarrow$ the dyadic Bell lift~\eqref{eq:supp:lift} is undefined. Only rung-$0$ distance-3 is recorded; out of scope.\\
\bottomrule
\end{tabular}
\end{table}



%%====================================================================
\section{Aggregate statistics}
\label{sec:agg}
% ============================================================

\subsection{Verdict breakdown}

\begin{center}
\begin{tabular}{@{}lr@{\hspace{3em}}r@{}}
\toprule
Verdict & Count & Share \\
\midrule
\vCONF{}\hspace{1em} all hypotheses hold and all rungs $d\ge 3$ & 30 & 69.8\,\% \\
\vCNTR{}\hspace{1em} hypotheses hold, some rung breaks $d\ge 3$ & \textbf{0}  & 0\,\% \\
\vINAP{}\hspace{1em} at least one hypothesis fails             & 9 & 20.9\,\% \\
\hspace{1em}\textit{of which conclusion still holds} & \textit{1} & \textit{2.3\,\%} \\
\vLIFT{}\hspace{1em} not SS-LP, lift undefined                 & 4  & 9.3\,\% \\
\midrule
Total tested & 43 & 100\,\% \\
\bottomrule
\end{tabular}
\end{center}

\subsection{Hypothesis pass rates}

Among the 39 Bell-liftable codes (excluding the 4 \vLIFT{} cases):
\begin{center}
\begin{tabular}{@{}lcc@{}}
\toprule
Hypothesis & Pass count & Pass rate \\
\midrule
\textbf{strong} ($\Sigma a = 2m-1$)        & 30/39 & 76.9\,\% \\
\textbf{bnd} ($\max a \le m-2$)            & 37/39 & 94.9\,\% \\
\textbf{(i)} unique bottom anchor           & 31/39 & 79.5\,\% \\
\textbf{(ii)} unique top anchor             & 30/39 & 76.9\,\% \\
\textbf{(iii)} seed + rung 1 distance-3     & 35/39 & 89.7\,\% \\
\addlinespace
\textbf{All five}                           & 30/39 & 76.9\,\% \\
\bottomrule
\end{tabular}
\end{center}

% ============================================================
%% ---- data section taken from the corrected si.tex ----
\section{Codes that confirm Theorem~4.4}
\label{sec:supp:confirms}
%%====================================================================

Table~\ref{tab:supp:confirms} lists the $30$ codes for which all five
hypotheses of Theorem~4.4 hold and every rung $r\in\{0,1,2,3,4,5\}$ preserves
distance $3$. The ``Worst viol.'' column is the maximum weight-$\le 2$ KL
violation across all six rungs, rounded to two significant figures.

\begin{longtable}{@{}lll cc l rr@{}}
\caption{The $30$ \vCONF{} codes.}\label{tab:supp:confirms}\\
\toprule
Code & §\,(in paper) & Group & $n_0$ & $m_0$ & $a$-vector & $\sum a$ & Worst viol. \\
\midrule
\endfirsthead
\multicolumn{8}{c}{\small\itshape Table~\ref{tab:supp:confirms} continued from previous page}\\
\toprule
Code & §\,(in paper) & Group & $n_0$ & $m_0$ & $a$-vector & $\sum a$ & Worst viol. \\
\midrule
\endhead
\midrule
\multicolumn{8}{r}{\small\itshape continued on next page}\\
\endfoot
\bottomrule
\endlastfoot
\texttt{steane}        & prototype & $2O$         & 7 & 4  & $(1,1,1,1,1,1,1)$        &  7 & \enot{3.1}{-16} \\
\texttt{c10}           & §4.1      & $C_{10}$     & 6 & 5  & $(1,1,1,1,2,3)$          &  9 & \enot{5.0}{-16} \\
\texttt{2i\_lambda34}  & §5.1      & $2I$         & 7 & 5  & $(1,1,1,1,1,2,2)$        &  9 & \enot{3.3}{-16} \\
\texttt{bd12\_type\_I} & App\,D.1  & $\mathrm{BD}_{12}$ & 7 & 6  & $(1,1,1,2,2,2,2)$ & 11 & \enot{4.2}{-16} \\
\texttt{bd16\_01}      & §6.3.1    & $\mathrm{BD}_{16}$ & 7 & 8  & $(1,2,2,2,2,3,3)$ & 15 & \enot{4.2}{-16} \\
\texttt{bd16\_02}      & §6.3.1    & $\mathrm{BD}_{16}$ & 7 & 8  & $(1,1,2,2,2,2,5)$ & 15 & \enot{2.8}{-16} \\
\texttt{bd16\_03}      & §6.3.1    & $\mathrm{BD}_{16}$ & 7 & 8  & $(0,1,2,2,3,3,4)$ & 15 & \enot{3.3}{-16} \\
\texttt{bd16\_04}      & §6.3.1    & $\mathrm{BD}_{16}$ & 7 & 8  & $(1,1,2,2,3,3,3)$ & 15 & \enot{4.3}{-16} \\
\texttt{bd16\_05}      & §6.3.1    & $\mathrm{BD}_{16}$ & 7 & 8  & $(1,1,2,2,2,3,4)$ & 15 & \enot{1.7}{-16} \\
\texttt{bd16\_06}      & §6.3.1    & $\mathrm{BD}_{16}$ & 7 & 8  & $(1,1,1,2,3,3,4)$ & 15 & \enot{4.3}{-16} \\
\texttt{bd16\_07}      & §6.3.1    & $\mathrm{BD}_{16}$ & 7 & 8  & $(1,1,1,2,2,4,4)$ & 15 & \enot{1.7}{-16} \\
\texttt{bd16\_08}      & §6.3.1    & $\mathrm{BD}_{16}$ & 7 & 8  & $(1,1,1,2,2,3,5)$ & 15 & \enot{4.0}{-16} \\
\texttt{bd16\_09}      & §6.3.1    & $\mathrm{BD}_{16}$ & 7 & 8  & $(1,1,1,1,3,4,4)$ & 15 & \enot{2.2}{-16} \\
\texttt{bd16\_10}      & §6.3.1    & $\mathrm{BD}_{16}$ & 7 & 8  & $(1,1,1,1,3,3,5)$ & 15 & \enot{2.8}{-16} \\
\texttt{bd18\_type\_I} & App\,D.4  & $\mathrm{BD}_{18}$ & 7 & 9  & $(2,2,2,2,3,3,3)$ & 17 & \enot{2.8}{-16} \\
\texttt{bd20}\rlap{$^\S$} & App\,D.5  & $\mathrm{BD}_{20}$ & 7 & 10 & $(1,2,2,3,3,4,4)$ & 19 & \enot{2.8}{-16} \\
\texttt{bd18\_type\_II}\rlap{$^\dagger$} & App\,D.4 & $\mathrm{BD}_{18}$ & 7 & 9  & $(1,1,2,3,3,3,4)$ & 17 & \enot{8.0}{-14} \\
\texttt{bd22}          & §D.6      & $\mathrm{BD}_{22}$ & 7 & 11 & $(1,2,3,3,3,4,5)$ & 21 & \enot{2.2}{-16} \\
\texttt{bd26}          & App\,D.8  & $\mathrm{BD}_{26}$ & 7 & 13 & $(2,2,3,3,4,5,6)$ & 25 & \enot{4.4}{-16} \\
\texttt{bd28}          & App\,D.9  & $\mathrm{BD}_{28}$ & 7 & 14 & $(2,3,3,4,4,5,6)$ & 27 & \enot{3.3}{-16} \\
\texttt{bd30}          & App\,D.10 & $\mathrm{BD}_{30}$ & 7 & 15 & $(2,2,3,3,4,7,8)$ & 29 & \enot{6.2}{-16} \\
\texttt{bd32}          & §6.3.2    & $\mathrm{BD}_{32}$ & 7 & 16 & $(2,3,4,4,5,6,7)$ & 31 & \enot{4.5}{-16} \\
\texttt{bd34}          & App\,D.11 & $\mathrm{BD}_{34}$ & 7 & 17 & $(2,2,3,4,5,8,9)$ & 33 & \enot{6.3}{-16} \\
\texttt{bd36}          & §6.3.3    & $\mathrm{BD}_{36}$ & 7 & 18 & $(2,3,4,5,6,7,8)$ & 35 & \enot{4.2}{-16} \\
\texttt{bd38}\rlap{$^\ddagger$} & App\,E.1  & $\mathrm{BD}_{38}$ & 8 & 19 & $(2,3,3,4,4,5,7,9)$ & 37 & \enot{1.6}{-15} \\
\texttt{bd72}          & App\,E.2  & $\mathrm{BD}_{72}$ & 8 & 36 & $(3,5,6,8,10,12,13,14)$ & 71 & \enot{3.2}{-16} \\
\texttt{bd74}          & App\,E.3  & $\mathrm{BD}_{74}$ & 8 & 37 & $(4,6,7,9,10,11,12,14)$ & 73 & \enot{2.5}{-16} \\
\texttt{bd76}          & App\,E.4  & $\mathrm{BD}_{76}$ & 8 & 38 & $(2,4,7,9,10,12,15,16)$ & 75 & \enot{2.8}{-16} \\
\texttt{bd78}          & App\,E.5  & $\mathrm{BD}_{78}$ & 8 & 39 & $(3,5,8,9,10,12,14,16)$ & 77 & \enot{2.2}{-16} \\
\texttt{bd80}          & App\,E.6  & $\mathrm{BD}_{80}$ & 8 & 40 & $(2,5,6,8,11,14,15,18)$ & 79 & \enot{2.9}{-16} \\
\end{longtable}

\noindent\footnotesize
$^\S$ \texttt{bd20}: our first transcription, taken from the typeset PDF, was wrong and failed rung-0 KL by $\enot{7.5}{-1}$; re-reading the codeword from the source \LaTeX{} of \texttt{arXiv:2504.20847} fixes it and every hypothesis then holds.
\normalsize


\noindent\footnotesize
$^\dagger$ \texttt{bd18\_type\_II}: PDF rendering merged the $1/18$ inside the
$|1100101\rangle$ coefficient; corrected to $\sqrt{b^2+d^2+1/18}$ and
remaining sign phases resolved numerically. Worst viol.\ is at the
$\sim 10^{-13}$ level, still $\sim 10^{6}$ times below tolerance $10^{-8}$.\\
$^\ddagger$ \texttt{bd38}: 16-variable constraint system solved by feasibility
LP; the solution reproduces the paper's stated identities
$\theta_1^2+\theta_2^2 = 1/38$ and $\theta_{12}^2 = 4/19$.
\normalsize

\paragraph{Remark on coverage.}
The BD$_{16}$ family is represented by $10$ distinct angle vectors, all of
which Confirm: Theorem~4.4 is not specific to a single representative but
applies to the whole family enumerated in the paper.

%%====================================================================
%% ---- data section taken from the corrected si.tex ----
\section{Codes for which Theorem~4.4 is inapplicable}
\label{sec:supp:inapp}
%%====================================================================

Table~\ref{tab:supp:inapp} lists the $9$ codes for which at least one
hypothesis of Theorem~4.4 fails. Red $\boldsymbol{\times}$ marks failed
hypotheses; ``Cncl?''\ records whether the conclusion (all $5$ rungs preserve
$d\ge 3$) still holds empirically.

\begin{longtable}{@{}lll cc cccccc c@{}}
\caption{The $9$ \vINAP{} codes.}\label{tab:supp:inapp}\\
\toprule
Code & §\,(paper) & Group & $n_0$ & $m_0$ & strong & bnd & (i) & (ii) & (iii) & Cncl? & Worst viol. \\
\midrule
\endfirsthead
\toprule
Code & §\,(paper) & Group & $n_0$ & $m_0$ & strong & bnd & (i) & (ii) & (iii) & Cncl? & Worst viol. \\
\midrule
\endhead
\bottomrule
\endlastfoot
\rowcolor{cSharpBg}
\texttt{2i\_perm\_inv} & §5.1      & $2I$         & 7 & 5  & \no  & \yes & \yes & \no  & \yes & \yes & \enot{2.8}{-16} \\
\texttt{2i\_lambda0}   & §5.1      & $2I$         & 7 & 5  & \no  & \yes & \no  & \no  & \no  & \no  & \enot{1.2}{0}  \\
\texttt{bd12\_type\_II}& App\,D.1  & $\mathrm{BD}_{12}$ & 7 & 6  & \no  & \yes & \no  & \no  & \no  & \no  & \enot{1.0}{0}  \\
\texttt{bd14}          & App\,D.2  & $\mathrm{BD}_{14}$ & 7 & 7  & \no  & \no  & \no  & \no  & \yes & \no  & \enot{7.1}{-2} \\
\texttt{bd16\_d3}      & App\,D.3  & $\mathrm{BD}_{16}$ & 7 & 8  & \no  & \yes & \no  & \no  & \no  & \no  & \enot{1.2}{0}  \\
\texttt{bd24}          & App\,D.7  & $\mathrm{BD}_{24}$ & 7 & 12 & \no  & \yes & \no  & \no  & \no  & \no  & \enot{1.7}{0}  \\
\texttt{bd64}          & §6.4.1    & $\mathrm{BD}_{64}$ & 8 & 32 & \no  & \yes & \no  & \no  & \yes & \no  & \enot{5.9}{-1} \\
\texttt{bd84}          & §6.4.2    & $\mathrm{BD}_{84}$ & 8 & 42 & \no  & \yes & \no  & \no  & \yes & \no  & \enot{4.5}{-1} \\
\texttt{2t\_cyclic}    & App\,C.1  & $2T$         & 7 & 2  & \no  & \no  & \no  & \no  & \no  & \no  & \enot{2.4}{-1} \\
\end{longtable}

\subsection{Sharpness: hypotheses are sufficient but not necessary}
\label{ssec:supp:sharpness}

The yellow-highlighted row \texttt{2i\_perm\_inv} (the $2I$
permutation-invariant 7-qubit code) is the only entry in
Table~\ref{tab:supp:inapp} where the conclusion of Theorem~4.4 holds despite
some hypothesis failing. Specifically $\Sigma a = 7 \ne 9 = 2m_0-1$, so
\textbf{strong} fails, and hypothesis (ii) is vacuously false because the
would-be top-anchor sum $2m_0-1=9$ is not achievable. Yet rungs $r=0,\ldots,5$
all preserve $d\ge 3$ at machine precision (worst viol.\
\enot{2.8}{-16}).

This single data point shows that Theorem~4.4's hypotheses are
\emph{sufficient but not necessary}. The theorem isolates a clean lift
certificate; \texttt{2i\_perm\_inv} suggests there is room for generalized
ladder theorems with weaker hypotheses, likely leveraging additional symmetry
such as permutation invariance.

\subsection{Failure-rung patterns}
\label{ssec:supp:failure-rungs}

Of the $9$ \vINAP{} codes, one (\texttt{2i\_perm\_inv}) keeps $d\ge3$ at every
rung anyway; for the remaining $8$ the conclusion fails at some rung. The rung
of first failure (Table~\ref{tab:supp:first-fail}) is consistent with which
hypothesis is violated:

\begin{table}[h]
\centering\small
\caption{First failing rung for the $8$ \vINAP{} codes whose conclusion also
fails. The same data appears in Table~\ref{tab:supp:rungs} as $r$-by-$r$
pass/fail marks.}\label{tab:supp:first-fail}
\begin{tabular}{@{}lcL{3.6cm}L{5.2cm}@{}}
\toprule
Code & First failing rung & Failing hypothesis (most relevant) & Comment \\
\midrule
\texttt{2i\_lambda0}   & rung 1 & strong, (i), (ii), (iii) & $\lambda^\star = 0$ representation \\
\texttt{bd12\_type\_II}& rung 1 & strong, (i), (ii), (iii) & $\Sigma a = 13 \ne 11$ \\
\texttt{bd16\_d3}      & rung 1 & strong, (i), (ii), (iii) & $X^5{\otimes}I^2$ family, not $X^7$-self-comp. \\
\texttt{bd24}          & rung 1 & strong, (i), (ii), (iii) & $\Sigma a = 25 \ne 23$ \\
\texttt{2t\_cyclic}    & rung 2 & strong, bnd, (i), (ii)   & $m_0 = 2$; transcription ambiguity \\
\texttt{bd64}          & rung 2 & strong, (i), (ii)        & $\Sigma a = 127 \ne 63$ \\
\texttt{bd84}          & rung 2 & strong, (i), (ii)        & $\Sigma a = 167 \ne 83$ \\
\texttt{bd14}          & rung 3 & strong, bnd, (i), (ii)   & $\Sigma a = 26,\ \max a = 6 > m_0{-}2=5$ \\
\bottomrule
\end{tabular}
\end{table}

Every one of the eight has a distance-3 seed: rung $0$ passes in all cases, so
no published codeword is in question here. What separates them is how far the
ladder carries before the missing hypothesis is felt.

\begin{itemize}[leftmargin=2em]
\item \emph{The first lift breaks} (rung $1$). \texttt{2i\_lambda0},
\texttt{bd12\_type\_II}, \texttt{bd16\_d3}, \texttt{bd24}: hypothesis (iii)
itself fails, so the first lifted rung already violates weight-$\le2$ KL and
the ladder never starts.
\item \emph{The strong sum is felt at rung $2$}. \texttt{2t\_cyclic},
\texttt{bd64}, \texttt{bd84}: hypothesis (iii) holds --- seed and first lift
both pass --- but $\Sigma a \ne 2m_0-1$, and rung $2$ ($m_r = 4m_0$) is
precisely where the carry-tag accounting first needs that identity to
propagate. \texttt{bd64} and \texttt{bd84} are large $\mathrm{BD}_{2m}$
families with $\Sigma a \approx 4m_0$.
\item \emph{One code survives a rung longer} (rung $3$). \texttt{bd14} has the
same hypothesis profile as \texttt{2t\_cyclic} --- (iii) holds, the strong-sum
and bounded-entry conditions both fail --- yet keeps $d\ge3$ through rung $2$.
The hypotheses classify which codes fail, not exactly when.
\end{itemize}

%%====================================================================
%% ---- data section taken from the corrected si.tex ----
\section{Codes outside the theorem's scope}
\label{sec:supp:nonliftable}
%%====================================================================

The four codes in Table~\ref{tab:supp:nonliftable} are not SS-LP (their
transversal gate is not diagonal in subset-sum sectors), so the dyadic Bell
lift~\eqref{eq:supp:lift} is undefined. We only report rung-$0$ distance-3
status as a baseline.

\begin{table}[h]
\centering\small
\caption{The $4$ \vLIFT{} codes (out of scope of Theorem~4.4).}
\label{tab:supp:nonliftable}
\begin{tabular}{@{}llccL{5.0cm}c@{}}
\toprule
Code & §\,(paper) & Group & $n_0,m_0$ & Reason out of scope & rung 0 $d\!\ge\!3$? \\
\midrule
\texttt{2o\_code}              & App\,C.2 & $2O$            & 7,\,2 & Non-diagonal transversal gate                         & \rOK \\
\texttt{c623\_from\_c523\_bd4} & App\,B.1 & $\mathrm{BD}_4$ & 6,\,2 & $((6,2,3))$ from $((5,2,3))$; not SS-LP               & \rOK \\
\texttt{c623\_from\_c523\_2t}  & App\,B.1 & $2T$            & 6,\,2 & $((6,2,3))$ from $((5,2,3))$; not SS-LP               & \rOK \\
\texttt{so5\_bd4}              & App\,B.2 & $\mathrm{BD}_4$ & 6,\,2 & SO(5) family; paper $W$ non-orthogonal as transcribed & \rFAIL \\
\bottomrule
\end{tabular}
\end{table}

\texttt{so5\_bd4}'s rung-$0$ failure ($\enot{2.6}{-1}$) is a transcription
issue: the paper's canonical $W$ matrix as rendered in the PDF has
$\mathrm{col}_0 \cdot \mathrm{col}_1 = -3 W_{15}W_{25} \ne 0$, i.e.\ it is not
orthogonal. Resolving this would require the source \LaTeX. Either way the
code is \vLIFT{}, so the failure is outside Theorem~4.4's scope.

%%====================================================================
%% ---- data section taken from the corrected si.tex ----
\section{Per-rung distance preservation}
\label{sec:supp:rungs}
%%====================================================================

Table~\ref{tab:supp:rungs} records, for every code, whether the weight-$\le 2$
KL condition holds at each rung $r\in\{0,1,2,3,4,5\}$. A \rOK{} cell means
pass within tolerance $10^{-8}$; \rFAIL{} means the worst weight-$\le 2$ KL
violation at that rung exceeds $10^{-8}$; \rNA{} means the lift is undefined
at that rung (this only occurs for \vLIFT{} codes, where only rung $0$ is
checked). Codes are grouped by verdict.

\begin{longtable}{@{}lll cccccc l@{}}
\caption{Per-rung weight-$\le 2$ KL pass/fail for every code.}
\label{tab:supp:rungs}\\
\toprule
Code & §\,(paper) & Verdict & $r_0$ & $r_1$ & $r_2$ & $r_3$ & $r_4$ & $r_5$ & Worst viol.\\
\midrule
\endfirsthead
\toprule
Code & §\,(paper) & Verdict & $r_0$ & $r_1$ & $r_2$ & $r_3$ & $r_4$ & $r_5$ & Worst viol.\\
\midrule
\endhead
\bottomrule
\endlastfoot
%% --- 29 CONFIRMS ---
\texttt{steane}        & prototype & \vCONF & \rOK & \rOK & \rOK & \rOK & \rOK & \rOK & \enot{3.1}{-16}\\
\texttt{c10}           & §4.1      & \vCONF & \rOK & \rOK & \rOK & \rOK & \rOK & \rOK & \enot{5.0}{-16}\\
\texttt{2i\_lambda34}  & §5.1      & \vCONF & \rOK & \rOK & \rOK & \rOK & \rOK & \rOK & \enot{3.3}{-16}\\
\texttt{bd12\_type\_I} & App\,D.1  & \vCONF & \rOK & \rOK & \rOK & \rOK & \rOK & \rOK & \enot{4.2}{-16}\\
\texttt{bd16\_01}      & §6.3.1    & \vCONF & \rOK & \rOK & \rOK & \rOK & \rOK & \rOK & \enot{4.2}{-16}\\
\texttt{bd16\_02}      & §6.3.1    & \vCONF & \rOK & \rOK & \rOK & \rOK & \rOK & \rOK & \enot{2.8}{-16}\\
\texttt{bd16\_03}      & §6.3.1    & \vCONF & \rOK & \rOK & \rOK & \rOK & \rOK & \rOK & \enot{3.3}{-16}\\
\texttt{bd16\_04}      & §6.3.1    & \vCONF & \rOK & \rOK & \rOK & \rOK & \rOK & \rOK & \enot{4.3}{-16}\\
\texttt{bd16\_05}      & §6.3.1    & \vCONF & \rOK & \rOK & \rOK & \rOK & \rOK & \rOK & \enot{1.7}{-16}\\
\texttt{bd16\_06}      & §6.3.1    & \vCONF & \rOK & \rOK & \rOK & \rOK & \rOK & \rOK & \enot{4.3}{-16}\\
\texttt{bd16\_07}      & §6.3.1    & \vCONF & \rOK & \rOK & \rOK & \rOK & \rOK & \rOK & \enot{1.7}{-16}\\
\texttt{bd16\_08}      & §6.3.1    & \vCONF & \rOK & \rOK & \rOK & \rOK & \rOK & \rOK & \enot{4.0}{-16}\\
\texttt{bd16\_09}      & §6.3.1    & \vCONF & \rOK & \rOK & \rOK & \rOK & \rOK & \rOK & \enot{2.2}{-16}\\
\texttt{bd16\_10}      & §6.3.1    & \vCONF & \rOK & \rOK & \rOK & \rOK & \rOK & \rOK & \enot{2.8}{-16}\\
\texttt{bd18\_type\_I} & App\,D.4  & \vCONF & \rOK & \rOK & \rOK & \rOK & \rOK & \rOK & \enot{2.8}{-16}\\
\texttt{bd18\_type\_II}& App\,D.4  & \vCONF & \rOK & \rOK & \rOK & \rOK & \rOK & \rOK & \enot{8.0}{-14}\\
\texttt{bd20}\rlap{$^\S$} & App\,D.5 & \vCONF & \rOK & \rOK & \rOK & \rOK & \rOK & \rOK & \enot{2.8}{-16}\\
\texttt{bd22}          & §D.6      & \vCONF & \rOK & \rOK & \rOK & \rOK & \rOK & \rOK & \enot{2.2}{-16}\\
\texttt{bd26}          & App\,D.8  & \vCONF & \rOK & \rOK & \rOK & \rOK & \rOK & \rOK & \enot{4.4}{-16}\\
\texttt{bd28}          & App\,D.9  & \vCONF & \rOK & \rOK & \rOK & \rOK & \rOK & \rOK & \enot{3.3}{-16}\\
\texttt{bd30}          & App\,D.10 & \vCONF & \rOK & \rOK & \rOK & \rOK & \rOK & \rOK & \enot{6.2}{-16}\\
\texttt{bd32}          & §6.3.2    & \vCONF & \rOK & \rOK & \rOK & \rOK & \rOK & \rOK & \enot{4.5}{-16}\\
\texttt{bd34}          & App\,D.11 & \vCONF & \rOK & \rOK & \rOK & \rOK & \rOK & \rOK & \enot{6.3}{-16}\\
\texttt{bd36}          & §6.3.3    & \vCONF & \rOK & \rOK & \rOK & \rOK & \rOK & \rOK & \enot{4.2}{-16}\\
\texttt{bd38}          & App\,E.1  & \vCONF & \rOK & \rOK & \rOK & \rOK & \rOK & \rOK & \enot{1.6}{-15}\\
\texttt{bd72}          & App\,E.2  & \vCONF & \rOK & \rOK & \rOK & \rOK & \rOK & \rOK & \enot{3.2}{-16}\\
\texttt{bd74}          & App\,E.3  & \vCONF & \rOK & \rOK & \rOK & \rOK & \rOK & \rOK & \enot{2.5}{-16}\\
\texttt{bd76}          & App\,E.4  & \vCONF & \rOK & \rOK & \rOK & \rOK & \rOK & \rOK & \enot{2.8}{-16}\\
\texttt{bd78}          & App\,E.5  & \vCONF & \rOK & \rOK & \rOK & \rOK & \rOK & \rOK & \enot{2.2}{-16}\\
\texttt{bd80}          & App\,E.6  & \vCONF & \rOK & \rOK & \rOK & \rOK & \rOK & \rOK & \enot{2.9}{-16}\\
\midrule
%% --- 10 INAPPLICABLE ---
\rowcolor{cSharpBg}
\texttt{2i\_perm\_inv} & §5.1      & \vINAP & \rOK & \rOK & \rOK & \rOK & \rOK & \rOK & \enot{2.8}{-16}\\
\texttt{2i\_lambda0}   & §5.1      & \vINAP & \rOK & \rFAIL & \rFAIL & \rFAIL & \rFAIL & \rFAIL & \enot{1.2}{0}\\
\texttt{bd12\_type\_II}& App\,D.1  & \vINAP & \rOK & \rFAIL & \rFAIL & \rFAIL & \rFAIL & \rFAIL & \enot{1.0}{0}\\
\texttt{bd14}          & App\,D.2  & \vINAP & \rOK & \rOK & \rOK & \rFAIL & \rFAIL & \rFAIL & \enot{7.1}{-2}\\
\texttt{bd16\_d3}      & App\,D.3  & \vINAP & \rOK & \rFAIL & \rFAIL & \rFAIL & \rFAIL & \rFAIL & \enot{1.2}{0}\\
\texttt{bd24}          & App\,D.7  & \vINAP & \rOK & \rFAIL & \rFAIL & \rFAIL & \rFAIL & \rFAIL & \enot{1.7}{0}\\
\texttt{bd64}          & §6.4.1    & \vINAP & \rOK & \rOK & \rFAIL & \rFAIL & \rFAIL & \rFAIL & \enot{5.9}{-1}\\
\texttt{bd84}          & §6.4.2    & \vINAP & \rOK & \rOK & \rFAIL & \rFAIL & \rFAIL & \rFAIL & \enot{4.5}{-1}\\
\texttt{2t\_cyclic}    & App\,C.1  & \vINAP & \rOK & \rOK & \rFAIL & \rFAIL & \rFAIL & \rFAIL & \enot{2.4}{-1}\\
\midrule
%% --- 4 NOT-LIFTABLE ---
\texttt{2o\_code}             & App\,C.2 & \vLIFT & \rOK   & \rNA & \rNA & \rNA & \rNA & \rNA & \enot{1.7}{-16}\\
\texttt{c623\_from\_c523\_bd4}& App\,B.1 & \vLIFT & \rOK   & \rNA & \rNA & \rNA & \rNA & \rNA & \enot{3.7}{-18}\\
\texttt{c623\_from\_c523\_2t} & App\,B.1 & \vLIFT & \rOK   & \rNA & \rNA & \rNA & \rNA & \rNA & \enot{3.9}{-17}\\
\texttt{so5\_bd4}             & App\,B.2 & \vLIFT & \rFAIL & \rNA & \rNA & \rNA & \rNA & \rNA & \enot{2.6}{-1}\\
\end{longtable}

% %%====================================================================
% \section{Notable cases and transcription notes}
% \label{sec:supp:notes}
% %%====================================================================

% \subsection{Theorem-sharpness families: \texttt{bd64} and \texttt{bd84}}

% These are the $\nu=5$ and $\nu=6$ BD-families in the paper's §6.4
% ($n_0=8$). Both rungs $0$ and $1$ satisfy weight-$\le 2$ KL, so the codewords
% are valid distance-$3$ SS-LP codes that even admit a first Bell lift. But
% $\Sigma a$ is approximately $2(2m_0{-}1)$ rather than $2m_0-1$, so the
% strong-sum identity does not propagate; at rung $2$ the carry-tag accounting
% fails and the lift breaks. This is precisely the case Theorem~4.4 rules out:
% the strong-sum identity (H1) is necessary for the lift's residue arithmetic
% to land in the right block beyond rung $1$.

% \subsection{Likely paper typos: \texttt{bd14}, \texttt{bd20}, \texttt{bd24}}

% \texttt{bd14}: $\Sigma a = 26 \ne 13 = 2m_0{-}1$ and $\max a = 6 > m_0{-}2 = 5$
% (both H1 and H2 fail).
% \texttt{bd20}: H1--H4 all pass, but the seed itself fails rung-$0$ KL by
% $\enot{7.5}{-1}$; the verifier's notes flag ``paper $Z$-gate list inconsistent
% with supports.''
% \texttt{bd24}: $\Sigma a = 25 \ne 23 = 2m_0{-}1$ (H1 fails).
% For these, the verifier records \vINAP{} with the side note that the
% transcription should be double-checked against the source paper's \LaTeX.

% \subsection{Numerically resolved phases (still confirms): \texttt{bd18\_type\_II}, \texttt{bd38}}

% \texttt{bd18\_type\_II}: the PDF rendering of the $|1100101\rangle$ coefficient
% merged the $1/18$ factor; corrected to $\sqrt{b^2+d^2+1/18}$ via norm
% constraints, and the remaining sign pattern resolved numerically. After this
% repair, all $6$ rungs pass; worst viol.\ \enot{8.0}{-14}.

% \texttt{bd38}: the $16$-variable constraint system ($8\times 16$ sign matrix
% $C$ plus the normalization $\Sigma = 1$) was solved by a feasibility LP. The
% solution reproduces the paper's stated identities
% $\theta_1^2 + \theta_2^2 = 1/38$ and $\theta_{12}^2 = 4/19$. All $6$ rungs
% pass; worst viol.\ \enot{1.6}{-15}.

% \subsection{The cyclic 2T family endpoints}

% The cyclic $2T$ family of App.\,C.1 has $m_0=2$, $\Sigma a = 7 \ne 3 =
% 2m_0-1$, so it is \vINAP{} regardless. Our reading of the paper's Eq.\,(18)
% recovers \texttt{steane} at the endpoint $\lambda^\star = 0$ and
% \texttt{2i\_perm\_inv} at $\lambda^\star = \sqrt 7$. Intermediate
% $\lambda^\star$ values fail rung $0$ in our transcription (likely an ambiguity
% in the paper's $c_1/c_3$ cyclic-orbit blocks); both endpoints are tested
% separately and behave as documented in their own rows of
% Tables~\ref{tab:supp:confirms} and~\ref{tab:supp:inapp}.

% \subsection{Coverage gap}

% The paper's §6.3.1 lists $12$ BD$_{16}$ angle vectors; $10$ are explicit enough
% to transcribe (Table~\ref{tab:supp:confirms}). The two not covered are
% $(1,1,1,2,2,2,6)$ and $(0,1,1,2,3,3,5)$, for which the paper gives no
% explicit $|0_L\rangle$ codewords. They are skipped from this verification.

% %%====================================================================
% \section{Reproducibility}
% \label{sec:supp:repro}
% %%====================================================================

% \paragraph{Pipeline.}
% \texttt{paper\_seeds\_2504\_20847.py} hard-codes each
% $(n_0, m_0, a, |0_L\rangle, |1_L\rangle)$ as a sparse JSON dictionary.
% \texttt{verify\_thm44\_paper.py} reads the seeds, iterates the Bell
% lift~\eqref{eq:supp:lift} for rungs $r = 0,\ldots,5$, enumerates all
% weight-$\le 2$ Paulis at each rung, computes the $2\times 2$ KL block,
% and records pass/fail (tolerance $10^{-8}$) plus worst KL violation.

% \paragraph{Numerical settings.}
% \begin{itemize}[leftmargin=2em]
% \item Precision: IEEE 754 FP64 throughout. The two symbolic-LP repairs
%   (\texttt{bd38}, \texttt{bd18\_type\_II}) use \texttt{gmpy2.mpq} exact
%   rationals during constraint solving.
% \item Tolerance: $\mathtt{tol} = 10^{-8}$. A weight-$\le 2$ KL violation $v$
%   at a rung counts as a fail iff $v > \mathtt{tol}$. All $29$ \vCONF{} cases
%   have $v \le \enot{8.0}{-14}$, six orders of magnitude inside tolerance.
% \item Max rungs: $r_{\max} = 5$. At rung $5$ a $7$-qubit seed reaches
%   $n_5 = 17$ qubits and $m_5 = 32 m_0$. The number of weight-$\le 2$ Pauli
%   strings at $n=17$ is $1 + 3\binom{17}{1} + 9\binom{17}{2} = 1276$, all
%   of which are enumerated.
% \end{itemize}

% \paragraph{Data files in this directory.}
% \begin{itemize}[leftmargin=2em]
% \item \texttt{thm44\_verification\_results.json} -- full per-code, per-rung
%   diagnostics (codewords, $2\times 2$ KL block, worst violation, witnesses).
% \item \texttt{thm44\_verification\_results.csv} -- one row per code, 24
%   columns; machine-friendly summary.
% \item \texttt{thm44\_verification\_summary.md} -- plain-text human-readable
%   report (superseded by this supplement).
% \end{itemize}

% %%====================================================================
% %% Closing notes (no section heading; flows into main paper text)
% %%====================================================================

% \paragraph{Concluding observations.}
% (i) Theorem~4.4 is empirically confirmed across the literature's explicit
% non-additive SS-LP codes: $29$ out of $29$ codes that satisfy its
% hypotheses preserve distance $3$ through $5$ dyadic Bell lifts at machine
% precision; zero counterexamples.
% (ii) Hypothesis sharpness is observed: codes failing strong-sum break at
% rung $2$ (\texttt{bd64}, \texttt{bd84}); codes failing (iii) break at
% rung $0$ or rung $1$, in a pattern consistent with which hypothesis is
% violated.
% (iii) Necessity is not implied: \texttt{2i\_perm\_inv} satisfies the
% conclusion despite failing two hypotheses, indicating that a generalized
% ladder theorem (perhaps leveraging permutation invariance) may admit weaker
% hypotheses.
% (iv) Coverage of the source paper is essentially complete; the two missing
% BD$_{16}$ angle vectors lack published $|0_L\rangle$ codewords.



%%====================================================================
%% End of supplement content
%%====================================================================

\section{Notable cases and transcription notes}
\label{sec:notes}
% ============================================================

\subsection{Theorem-sharpness families: \texttt{bd64} and \texttt{bd84}}

These two are the $\nu=5$ and $\nu=6$ BD-families in the paper's Sec.\,6.4
($n_0=8$). They are interesting because rungs 0 and 1 \emph{do} satisfy
weight-$\le2$ KL: the codewords are valid distance-3 SS-LP codes and even
admit a first Bell lift. But $\Sigma a$ is approximately $2\cdot(2m{-}1)$
rather than $2m{-}1$, so the strong-sum identity does not propagate. At
rung 2 the carry-tag accounting fails and the lift breaks.

This is the closest the dataset comes to ``looking like a counterexample.''
It is in fact \emph{precisely the case the theorem rules out}: the strong-sum
identity (H1) is necessary for the lift's residue arithmetic to land in the
right block beyond rung 1.

\subsection{Angle vectors outside the hypotheses: \texttt{bd14}, \texttt{bd24}}

Neither code is defective; their angle vectors simply do not meet the
arithmetic conditions the theorem asks for, so Theorem~4.4 makes no claim
about them.

\begin{itemize}[leftmargin=2em]
\item \texttt{bd14}: $\Sigma a = 26 \ne 13 = 2m{-}1$ and
$\max a = 6 > m{-}2 = 5$; H1 and H2 both fail.
\item \texttt{bd24}: $\Sigma a = 25 \ne 23 = 2m{-}1$; H1 fails. The seed
itself \emph{is} a distance-3 code --- only the strong-sum identity is
violated.
\end{itemize}

\subsection{A transcription error of ours: \texttt{bd20}}

\texttt{bd20} was initially read off the typeset PDF, and that reading failed
rung-0 KL by $\enot{7.5}{-1}$ --- large enough that the fault was clearly in
the codeword, not in the theorem. Re-reading the codeword from the source
\LaTeX{} of \texttt{arXiv:2504.20847} resolves it: with the corrected
amplitudes, $\Sigma a = 19 = 2m{-}1$, every hypothesis holds, and all six
rungs pass with worst violation $\enot{2.8}{-16}$. \texttt{bd20} is counted
among the \vCONF{} codes.

We record this because it is the one place where the pipeline initially
disagreed with the theorem, and the disagreement was ours.

\subsection{Numerically resolved phases (still confirms):
\texttt{bd18\_type\_II}, \texttt{bd38}}

\begin{itemize}[leftmargin=2em]
\item \texttt{bd18\_type\_II}: the PDF rendering of the $|1100101\rangle$
coefficient merged the $1/18$ factor; corrected to
$\sqrt{b^2+d^2+1/18}$ via norm constraints. The sign pattern of the other
terms was resolved numerically (the PDF $\pm$ rendering was unreliable).
After this repair, all 6 rungs pass; worst viol.\ \enot{8.0}{-14}.
\item \texttt{bd38}: the 16-variable constraint system (an $8\times16$ sign
matrix $C$ plus $\Sigma=1$) was solved by a feasibility LP. The solution
reproduces the paper's stated identities $\theta_1^2 + \theta_2^2 = 1/38$ and
$\theta_{12}^2 = 4/19$.  All 6 rungs pass; worst viol.\ \enot{1.6}{-15}.
\end{itemize}

\subsection{Coverage gap: two BD$_{16}$ angle vectors without published codewords}

The paper's \S 6.3.1 lists 12 BD$_{16}$ angle vectors; 10 of them are explicit
enough to transcribe. The two not covered are
$(1,1,1,2,2,2,6)$ and $(0,1,1,2,3,3,5)$, for which the paper gives no
explicit $|0_L\rangle$ codewords. They are skipped from this verification.

\subsection{The \texttt{2t\_cyclic} family endpoints}

The cyclic $2T$ family of App.\,C.1 has $m_0=2$, $\Sigma a = 7 \ne 3 = 2m{-}1$
(H1 fails); it is \vINAP{} regardless. Our Eq.\,(18) reading recovers
\texttt{steane} at the endpoint $\lambda^\star = 0$ and \texttt{2i\_perm\_inv}
at $\lambda^\star = \sqrt7$. Intermediate $\lambda^\star$ values fail rung 0
in our transcription (likely an ambiguity in the paper's $c_1$/$c_3$
cyclic-orbit block), but both endpoints are tested separately and behave as
documented in their own rows of Table~\ref{tab:supp:confirms}.

% ============================================================
\section{Reproducibility}
\label{sec:repro}
% ============================================================

\subsection{Pipeline}

\begin{center}
\begin{tabular}{@{}rl@{}}
\toprule
Step & Action \\
\midrule
1 & \texttt{paper\_seeds\_2504\_20847.py} hard-codes each $(n,m,a,|0_L\rangle,|1_L\rangle)$\\
2 & \texttt{verify\_thm44\_paper.py} reads the seeds and:\\
  & \quad (a) iterates the Bell lift~\eqref{eq:supp:lift} for rungs $r=0,\ldots,5$;\\
  & \quad (b) at each rung, enumerates all weight-$\le 2$ Paulis on $n_r$ qubits;\\
  & \quad (c) computes the $2\times 2$ KL block and checks the scalar condition;\\
  & \quad (d) records pass/fail (tolerance $10^{-8}$) and worst KL violation;\\
3 & Writes \texttt{thm44\_verification\_results.\{json,csv\}} and the
       human-readable summary.\\
\bottomrule
\end{tabular}
\end{center}

\subsection{Output files (this directory)}

\begin{center}
\begin{tabular}{@{}lL{9cm}@{}}
\toprule
File & Content \\
\midrule
\texttt{thm44\_verification\_results.json} & Full per-code, per-rung
diagnostics (codewords, KL block, worst violations, witnesses).\\
\texttt{thm44\_verification\_results.csv}  & One row per code, 24 columns;
machine-friendly summary.\\
\texttt{thm44\_verification\_summary.md}   & Plain-text human report
(superseded by this LaTeX version).\\
\texttt{thm44\_verification\_report.tex}   & This report.\\
\bottomrule
\end{tabular}
\end{center}

\subsection{Numerical settings}

\begin{itemize}[leftmargin=2em]
\item Precision: IEEE 754 FP64 throughout. Symbolic LP repairs
(\texttt{bd38}, \texttt{bd18\_type\_II}) use \texttt{gmpy2.mpq} exact rationals.
\item Tolerance: $\mathtt{tol} = 10^{-8}$. A weight-$\le 2$ KL violation $v$ at
a rung counts as a fail iff $v > \mathtt{tol}$. All \vCONF{} cases have
$v \le \enot{8.0}{-14}$, twelve orders of magnitude inside tolerance.
\item Max rungs: $r_{\max} = 5$. At rung 5 a $7$-qubit seed reaches
$n_5 = 17$ qubits and $m_5 = 2^5 m_0$. The number of weight-$\le 2$ Paulis at
$n=17$ is $1 + 3\binom{17}{1} + 9\binom{17}{2} = 1 + 51 + 1224 = 1276$, all of
which are enumerated.
\end{itemize}

% ============================================================
\section{Conclusions}
\label{sec:conclusions}
% ============================================================

\begin{enumerate}[leftmargin=2em]
\item \textbf{Theorem 4.4 is empirically confirmed across the literature's
explicit non-additive SS-LP codes.} 30 out of 30 codes that satisfy its
hypotheses preserve distance 3 through 5 dyadic Bell lifts at machine
precision. Zero counterexamples.
\item \textbf{Hypothesis sharpness is observed.} The codes that fail at least
one hypothesis split predictably: those that fail strong-sum break at
rung 2 (\texttt{bd64}, \texttt{bd84}); those that fail (iii) break at
rung 0 or rung 1.
\item \textbf{Necessity is not implied.} \texttt{2i\_perm\_inv} satisfies the
conclusion despite failing two hypotheses, indicating room for generalized
ladder theorems beyond Theorem 4.4's scope (probably leveraging additional
symmetry such as permutation invariance).
\item \textbf{Coverage of the source paper is essentially complete.} 43 of the
$\sim$45 explicit non-additive codes in arXiv:2504.20847 are covered; the
two missing BD$_{16}$ angle vectors lack published $|0_L\rangle$ codewords.
\end{enumerate}

\paragraph{Acknowledgement of dependence on transcription.}
Several codes' coefficients were transcribed from a PDF source, and three
cases needed care: in \texttt{bd18\_type\_II} and \texttt{bd38} the
coefficients were recovered via norm constraints and a feasibility LP, and in
\texttt{bd20} the PDF reading was simply wrong and was replaced by the source
\LaTeX. All three then satisfy every hypothesis. \texttt{bd14} and
\texttt{bd24} are a different matter: their angle vectors fall outside the
theorem's arithmetic conditions, which is a property of the codes and not a
transcription problem. A single-source-of-truth implementation $-$ generating
the codeword tables from the same Python file that drives the verification $-$
would remove the transcription risk entirely.

\end{document}
